
\definecolor{brown}{rgb}{0.65,0.16,0.16}
\definecolor{orange}{rgb}{1,0.65,0}
\definecolor{fuchsia}{rgb}{1,0,1}
\definecolor{gray}{rgb}{0.50,0.50,0.50}
\definecolor{purple}{rgb}{0.50,0,0.50}
\definecolor{violet}{rgb}{0.93,0.51,0.93}
\definecolor{maroon}{rgb}{0.50,0,0}
\Problem{Servidors}

\Statement

Volem simular el comportament d'una instal·lació d'$n$ servidors. Cada servidor està identificat amb un codi alfa-numèric i podem fer les instruccions següents sobre ells:

\begin{itemize}
\item \texttt{ADD server\_id job\_value}: assigna un nou treball amb valor \texttt{job\_value} (que sempre serà un número natural) al servidor \texttt{server\_id}.
\item \texttt{EXECUTE server\_id q}: executa els \texttt{q} treballs amb més valor del servidor \texttt{server\_id}. El valor \texttt{q} sempre és no negatiu.
\item \texttt{TRANSFER from\_id to\_id}: transfereix el treball més valuós del servidor amb identificador \texttt{from\_id} al servidor amb identificador \texttt{to\_id}. Si el servidor \texttt{from\_id} no té cap treball pendent, s'ignora la instrucció.
\end{itemize}

Heu de tenir en compte que de vegades s'utilitzen identificadors de servidors que no existeixen. Si és el cas, s'ha d'indicar convenientment.


\Input

L'entrada comença amb un natural $n$ seguit de l'identificador d'$n$ servidors. A continuació segueix una seqüència d'instruccions.

\Output

Per cada instrucció \texttt{ADD}, \texttt{EXECUTE} i \texttt{TRANSFER} s'ha d'escriure el valor del treball pendent més valuós de la/es cua/es implicada/es un cop realitzada la instrucció (escriu "-" en cas que el servidor no en tingui de pendents) o "identificador incorrecte" en cas que algun identificador de servidor ho sigui.

\medskip

Al final cal escriure el treball executat més valuós (o $-1$ si no s'ha executat cap treball) i un llistat amb el valor del treball més valuós pendent d'executar en cadascun dels servidors. Aquest llistat està ordenat de forma creixent segons l'identificador dels servidors. Quan un servidor no té cap treball pendent, s'indica amb "-". 

\medskip

Seguiu el format dels exemples.

\Sample

